\mxunit{automatismes-ch06}{Calcul algébrique, équations et inéquations}

\begin{fluencysheet}{jeu.~14~janv.~2027}{Puissances entières relatives}
\mxfq{Échauffement · Flash}{PUI}{Calculer $2^{-3}$, $10^{-2}$ et $(-3)^2$.}{$\dfrac18$\,; $0{,}01$\,; $9$.}{}
\mxfq{Cœur · Méthode}{PUI}{Simplifier $\dfrac{2^3\times2^{-5}}{2^{-1}}$.}{$\dfrac{2^{-2}}{2^{-1}}=2^{-2+1}=2^{-1}=\dfrac12$.}{}
\mxfq{Cœur + · Méthode}{PUI}{Pour $x\neq0$, écrire $(x^2)^{-3}\times x^7$ sous la forme $x^n$.}{$x^{-6}\times x^7=x^{1}=x$.}{}
\mxfq{Retour · Flash}{STA}{Médiane de la série $4,\,7,\,9,\,10,\,12,\,15$\,?}{Effectif pair~: $\dfrac{9+10}{2}=9{,}5$.}{}
\mxfq{Fil rouge · Reconnaître}{SOL}{Quel solide a une base en forme de disque, une surface latérale courbe et un unique sommet\,?}{Un cône de révolution.}{}
\end{fluencysheet}

\begin{fluencysheet}{lun.~18~janv.~2027}{Racines carrées\,; $\sqrt{a^2}=|a|$}
\mxfq{Échauffement · Flash}{CAL}{Calculer $\sqrt{36}$ et $\sqrt{0{,}09}$.}{$6$\,; $0{,}3$.}{}
\mxfq{Cœur · Méthode}{CAL}{Écrire $\sqrt{72}$ sous la forme $a\sqrt2$ et calculer $\sqrt{(-5)^2}$.}{$\sqrt{72}=\sqrt{36\times2}=6\sqrt2$\,; $\sqrt{(-5)^2}=|{-}5|=5$.}{}
\mxfq{Cœur + · Vrai/faux}{CAL}{Vrai ou faux\,? Justifier. $\sqrt{9+16}=\sqrt9+\sqrt{16}$.}{Faux. $\sqrt{25}=5$ mais $3+4=7$. Seul le produit se «~sépare~»~: $\sqrt{ab}=\sqrt a\sqrt b$.}{}
\mxfq{Retour · Flash}{PUI}{Calculer $(3^2)^3$.}{$3^{2\times3}=3^6=729$.}{}
\mxfq{Fil rouge · Vrai/faux}{DEM}{Vrai ou faux\,? Justifier. Les diagonales d'un losange sont perpendiculaires.}{Vrai. Propriété du losange (ses diagonales sont aussi médiatrices l'une de l'autre).}{}
\end{fluencysheet}

\begin{fluencysheet}{mar.~19~janv.~2027}{Expressions fractionnaires\,; formes développée, réduite ou factorisée}
\mxfq{Échauffement · Flash}{CAL}{Calculer $\dfrac12+\dfrac13$.}{$\dfrac56$.}{}
\mxfq{Cœur · Méthode}{LIT}{Écrire sous une seule fraction $\dfrac1x+\dfrac1{x+1}$ (pour $x\neq0$ et $x\neq-1$).}{$\dfrac{x+1+x}{x(x+1)}=\dfrac{2x+1}{x(x+1)}$.}{}
\mxfq{Cœur + · Méthode}{DEV}{Factoriser $x^2-4$ puis simplifier $\dfrac{x^2-4}{x-2}$ pour $x\neq2$.}{$x^2-4=(x-2)(x+2)$\,; donc $\dfrac{x^2-4}{x-2}=x+2$.}{Le choix de la forme dépend de l'objectif~: factorisée pour simplifier, développée pour comparer.}
\mxfq{Retour · Flash}{SUB}{Calculer $2a-3b$ pour $a=2$ et $b=-1$.}{$4+3=7$.}{}
\mxfq{Fil rouge · Flash}{REP}{$ABCDEFGH$ est un pavé droit avec $AB=4$, $AD=3$, $AE=2$. Dans le repère $(A;B,D,E)$, quelles sont les coordonnées de $G$\,?}{$G(4;3;2)$~: $G$ est le sommet opposé à $A$.}{Une troisième coordonnée apparaît dans l'espace.}
\end{fluencysheet}

\begin{fluencysheet}{jeu.~21~janv.~2027}{Inégalités~: somme, multiplication par un réel}
\mxfq{Échauffement · Flash}{CMP}{Compléter par $<$ ou $>$~: $-3\ \ldots\ -2$\,; $\dfrac13\ \ldots\ 0{,}3$.}{$-3<-2$\,; $\dfrac13>0{,}3$ ($0{,}333\ldots>0{,}3$).}{}
\mxfq{Cœur · Méthode}{EQ}{On a $a<b$. Comparer $-2a$ et $-2b$, puis $a+5$ et $b+5$.}{$-2a>-2b$ (multiplier par un négatif renverse l'inégalité)\,; $a+5<b+5$.}{}
\mxfq{Cœur + · Méthode}{EQ}{Sachant $2<x<3$ et $1<y<4$, encadrer $x+y$ puis $x-y$.}{$3<x+y<7$\,; pour $x-y$~: $-y\in]-4;-1[$ donc $-2<x-y<2$.}{Piège~: soustraire les inégalités membre à membre.}
\mxfq{Retour · Flash}{CAL}{Calculer $(-3)\times(-4)+2\times(-5)$.}{$12-10=2$.}{}
\mxfq{Fil rouge · Problème}{DIV}{On découpe une plaque de $120$ cm sur $84$ cm en carrés identiques, sans perte. Quel est le plus grand côté possible\,?}{$\pgcd(120\,;\,84)=12$~: carrés de $12$ cm de côté ($10\times7=70$ carrés).}{}
\end{fluencysheet}

\begin{fluencysheet}{lun.~25~janv.~2027}{Comparaison par différence et par quotient}
\mxfq{Échauffement · Flash}{CMP}{Comparer $\dfrac23$ et $\dfrac34$ en étudiant le signe de la différence.}{$\dfrac23-\dfrac34=-\dfrac1{12}<0$~: $\dfrac23<\dfrac34$.}{}
\mxfq{Cœur · Méthode}{LIT}{Comparer $x^2+1$ et $2x$ pour tout réel $x$.}{$x^2+1-2x=(x-1)^2\geqslant0$~: $x^2+1\geqslant2x$, égalité si et seulement si $x=1$.}{}
\mxfq{Cœur + · Méthode}{CMP}{Comparer $\dfrac75$ et $\dfrac86$ par quotient (nombres positifs).}{$\dfrac{7/5}{8/6}=\dfrac{42}{40}>1$, donc $\dfrac75>\dfrac86$.}{}
\mxfq{Retour · Tableau}{LEC}{Précipitations mensuelles (données fictives).\par\smallskip\mxarraystretch{1.15}\begin{tabular}{|c|c|c|c|c|}\hline
Mois & Janv. & Fév. & Mars & Avr. \\ \hline
Pluie (mm) & 280 & 300 & 240 & 120 \\ \hline
\end{tabular}\par Quel mois est le plus pluvieux\,? Quel est le cumul sur ces $4$ mois\,?}{Février ($300$ mm)\,; $940$ mm.}{}
\mxfq{Fil rouge · Tableau → formule}{REG}{$x$~: $-2,\,-1,\,0,\,1$ et $f(x)$~: $4,\,1,\,0,\,1$. La fonction est-elle affine\,? Proposer une formule.}{Non (variations de pas non constants). $f(x)=x^2$ convient.}{Test~: $f(1)-f(0)=1$ mais $f(0)-f(-1)=-1$.}
\end{fluencysheet}

\begin{fluencysheet}{mar.~26~janv.~2027}{Équations et inéquations du 1\textsuperscript{er} degré\,; $x^2=a$\,; produit nul}
\mxfq{Échauffement · Flash}{EQ}{Résoudre $3x-5=7$.}{$x=4$.}{}
\mxfq{Cœur · Méthode}{EQ}{Résoudre $2x+3\leqslant7-x$.}{$3x\leqslant4$~: $x\leqslant\dfrac43$\,; $S=\left]-\infty;\dfrac43\right]$.}{}
\mxfq{Cœur + · Méthode}{EQ}{Résoudre $(x-2)(2x+6)=0$ puis $x^2=7$.}{Produit nul~: $x=2$ ou $x=-3$\,; $x^2=7$~: $x=\sqrt7$ ou $x=-\sqrt7$.}{}
\mxfq{Retour · Flash}{CMP}{Comparer $1{,}4$ et $\sqrt2$ sans calculatrice.}{$1{,}4^2=1{,}96<2$ donc $1{,}4<\sqrt2$ (la fonction carré est croissante sur $[0;+\infty[$).}{}
\mxfq{Fil rouge · Flash}{GEO}{$ABC$ est isocèle en $A$ avec $\widehat{BAC}=40\degre$. Quelle est la mesure de $\widehat{ABC}$\,?}{$(180-40)\div2=70\degre$.}{}
\end{fluencysheet}

\begin{fluencysheet}{jeu.~28~janv.~2027}{Équation quotient $\frac{A(x)}{B(x)}=k$}
\mxfq{Échauffement · Flash}{EQ}{Pour quelle valeur de $x$ l'expression $\dfrac{1}{x-3}$ n'est-elle pas définie\,?}{$x=3$.}{}
\mxfq{Cœur · Méthode}{EQ}{Résoudre $\dfrac{2x-1}{x+3}=0$.}{$x\neq-3$ et $2x-1=0$~: $x=\dfrac12$.}{}
\mxfq{Cœur + · Méthode}{EQ}{Résoudre $\dfrac{x+1}{x-2}=3$.}{$x\neq2$\,; $x+1=3(x-2)$, soit $x=\dfrac72$. Contrôle~: $\dfrac{4{,}5}{1{,}5}=3$.}{}
\mxfq{Retour · Développer}{DEV}{Développer et réduire $(x+2)(x-5)$.}{$x^2-3x-10$.}{}
\mxfq{Fil rouge · Lecture de code}{ALG}{Que renvoient $g(3)$ et $g(-5)$\,?\par\smallskip{\ttfamily\small \hspace*{0.0em}def g(x):\\ \hspace*{2.4em}if x \textgreater{} 0:\\ \hspace*{4.8em}return 2*x\\ \hspace*{2.4em}else:\\ \hspace*{4.8em}return -x\par}}{$g(3)=6$ et $g(-5)=5$.}{Condition évaluée avant le calcul~: la branche dépend du signe de $x$.}
\end{fluencysheet}

\begin{fluencysheet}{lun.~1\textsuperscript{er}~févr.~2027}{Signe d'un produit ou d'un quotient\,; tableaux de signes}
\mxfq{Échauffement · Flash}{EQ}{Signe de $2x-4$\,?}{Négatif pour $x<2$, nul en $2$, positif pour $x>2$.}{}
\mxfq{Cœur · Méthode}{EQ}{Dresser le tableau de signes de $(x-1)(x+3)$.}{Racines $-3$ et $1$~: positif sur $]-\infty;-3[$, négatif sur $]-3;1[$, positif sur $]1;+\infty[$.}{Signe du produit~: règle des signes colonne par colonne.}
\mxfq{Cœur + · Méthode}{EQ}{Résoudre $\dfrac{x+2}{x-1}\geqslant0$.}{Numérateur positif pour $x\geqslant-2$, dénominateur positif pour $x>1$\,; $S=]-\infty;-2]\cup]1;+\infty[$ ($1$ est exclu).}{}
\mxfq{Retour · Flash}{PRB}{Jeu de $32$ cartes~: probabilité de tirer un cœur\,? un roi\,?}{$\dfrac{8}{32}=\dfrac14$\,; $\dfrac{4}{32}=\dfrac18$.}{}
\mxfq{Fil rouge · Réfuter}{LIT}{Réfuter~: «~pour tout réel $x$, $x^2\geqslant x$~».}{Contre-exemple~: $x=\dfrac12$~: $x^2=\dfrac14<\dfrac12=x$.}{Un seul contre-exemple suffit à réfuter une proposition universelle.}
\end{fluencysheet}
